\subsection{Compound Interest} 
One common place to find an exponential model is \emph{compound interest}, in particular, when we start to charge interest \emph{on the interest}.
%\begin{note}
%``Interest'' is a fee paid by a borrower or a lender. We will refer to ``loans,'' ``deposits,'' ``investments,'' or ``funds,'' interchangeably. Any situation where one party gives money to another with the intention of getting more money back later can be thought of in terms of interest.
%\end{note}
%What happens if we charge interest \emph{on the interest}?
\begin{example}
Suppose you borrow \$1000 at 10\% interest. What happens each year?
\begin{cpic}{}{5}
\regtext{-8}{-2}{Year};
\regtextml{-5}{-2}{1}{Amount\\owed at start};
\regtext{1}{-2}{Interest};
\regtextml{7}{-2}{1}{Amount\\owed at end};
\foreach \x in {-7,-3,5} \draw (\x,0) -- (\x,-5);
\draw (-9,-2) -- (9,-2);
\foreach \year in {1,2,3} \regtext{-8}{-2-\year}{$\year$};
\end{cpic}
\begin{itemize}
\item Each year, we multiply by \makeblank
\item This is an exponential model: $P(t)=\makeblank$
\end{itemize}
\end{example}
We could also charge interest more than once a year.
\begin{example}
What happens if we borrow \$1000 at 10\% interest, but we ``compound'' the interest every six months?
\begin{cpic}{}{6}
\regtext{-8}{-2}{Year};
%\regtext{-3}{-1}{Amount owed at start};
\regtextml{-4}{-2}{1}{Amount\\owed at start};
\regtext{2}{-2}{Interest};
%\regtext{7}{-1}{Amount owed at end};
\regtextml{8}{-2}{1}{Amount\\owed at end};
\foreach \x in {-6,-2,6} \draw (\x,0) -- (\x,-6);
\draw (-10,-2) -- (10,-2);
\foreach \year in {1,2} {
	\draw (-9,-\year*2-1) node[anchor=east]{\Large\year};
	\regtextleft{-9}{-\year*2-1}{1st half};
	\regtextleft{-9}{-\year*2-2}{2nd half};
}
\end{cpic}
\begin{itemize}
\item Now we multiply by \makeblank[1in], but we do it \makeblank[1in] times!
\item This is another exponential model: $P(t)=\makeblank$
\end{itemize}
\end{example}
\begin{definition}[Compound Interest]
If amount $P_0$ is invested at an annual interest rate $r$, compounded $n$ times per year, then the balance after $t$ years is given by
\[
P(t)=P_0\left(1+\frac{r}{n}\right)^{nt}
\]
\end{definition}